\documentclass[aspectratio=169,11pt]{beamer}

\usetheme{Madrid}
\usecolortheme{whale}
\setbeamertemplate{navigation symbols}{}
\setbeamertemplate{footline}[frame number]

\usepackage[utf8]{inputenc}
\usepackage[T1]{fontenc}
\usepackage{booktabs}
\usepackage{xcolor}
\usepackage{hyperref}
\usepackage{listings}

\lstset{
    basicstyle=\ttfamily\scriptsize,
    breaklines=true,
    frame=single,
    backgroundcolor=\color{gray!10},
    keywordstyle=\color{blue!70},
    commentstyle=\color{green!50!black},
    stringstyle=\color{red!60},
    showstringspaces=false,
    tabsize=2
}

\definecolor{stepbg}{RGB}{232,245,233}
\definecolor{stepborder}{RGB}{76,175,80}
\definecolor{warnbg}{RGB}{255,243,224}
\definecolor{warnborder}{RGB}{255,152,0}
\definecolor{aibg}{RGB}{243,229,245}
\definecolor{aiborder}{RGB}{156,39,176}

\title[EECE5155 -- Hands-On 1]{ESP32 Firmware on Wokwi Simulator}
\subtitle{Hands-On Session 1: Write Code Before Hardware Arrives}
\author{Eduardo Baena}
\date{Monday, March 23, 2026}
\institute{Northeastern University\\EECE5155: Wireless Sensor Networks and IoT}

\begin{document}

\begin{frame}
    \titlepage
    \vspace{-0.8cm}
    \begin{center}
        {\small \url{https://wokwi.com/}}
    \end{center}
\end{frame}

%==============================================================================
\section{What is Wokwi?}
%==============================================================================

\begin{frame}{What is Wokwi?}
    \textbf{Wokwi} is a free, browser-based electronics simulator.

    \vspace{0.2cm}
    \begin{columns}
        \column{0.55\textwidth}
        \textbf{Key facts:}
        \begin{itemize}
            \item Runs entirely in the browser
            \item Free to use; account only required to save projects
            \item Supports ESP32, ESP32-S2, ESP32-S3, ESP32-C3, Arduino, STM32, RPi Pico
            \item Real WiFi simulation (connects to actual internet)
            \item MQTT, HTTP, WebSocket all work from the simulator
            \item Arduino framework + ESP-IDF supported
            \item Free for personal and educational use
        \end{itemize}

        \column{0.45\textwidth}
        \fcolorbox{stepborder}{stepbg}{
        \begin{minipage}{0.9\textwidth}
            \vspace{0.1cm}
            \textbf{Wokwi vs.\ alternatives}
            \small
            \begin{itemize}
                \item Tinkercad: only Arduino, lacks WiFi/MQTT
                \item Proteus: paid license, heavy install
                \item Wokwi: free, instant, full network stack
            \end{itemize}
            \vspace{0.1cm}
        \end{minipage}
        }

        \vspace{0.2cm}
        \begin{block}{URL}
            \centering
            \url{https://wokwi.com/}
        \end{block}
    \end{columns}
\end{frame}

\begin{frame}{Wokwi Interface Tour}
    \textbf{When you open a project, you see four areas:}

    \vspace{0.2cm}
    \begin{enumerate}
        \item \textbf{Code Editor (left panel)}
        \begin{itemize}
            \item Your \texttt{.ino} sketch (Arduino C++)
            \item Tabs for additional files (\texttt{diagram.json}, \texttt{libraries.txt})
            \item Syntax highlighting, basic autocomplete
        \end{itemize}

        \item \textbf{Simulation View (right panel)}
        \begin{itemize}
            \item Virtual breadboard with your components
            \item Click components to interact (change sensor values, press buttons)
            \item Drag to reposition, scroll to zoom
        \end{itemize}

        \item \textbf{Serial Monitor (bottom panel)}
        \begin{itemize}
            \item Appears when simulation runs
            \item Shows \texttt{Serial.print()} output
            \item Has input field for sending data to your sketch
        \end{itemize}

        \item \textbf{Toolbar (top)}
        \begin{itemize}
            \item \textbf{Green Play}: start simulation
            \item \textbf{Red Stop}: stop simulation
            \item \textbf{Save / Share / Download}
        \end{itemize}
    \end{enumerate}
\end{frame}

\begin{frame}{Adding Components and Wiring}
    \begin{columns}
        \column{0.5\textwidth}
        \textbf{How to add a component:}
        \begin{enumerate}
            \item Click the \textbf{blue ``+'' button} on the simulation view
            \item Search by name (e.g., ``DHT22'', ``LED'', ``servo'')
            \item Click the component to add it
            \item Drag it into position on the board
        \end{enumerate}

        \vspace{0.2cm}
        \textbf{How to wire:}
        \begin{enumerate}
            \item Click on a component pin (a green dot appears)
            \item Click on the destination pin (ESP32 GPIO, GND, 3V3)
            \item A wire is drawn automatically
            \item Click a wire to delete it
        \end{enumerate}

        \column{0.5\textwidth}
        \textbf{The \texttt{diagram.json} file:}

        \vspace{0.1cm}
        Every component and wire is stored in \texttt{diagram.json}. You can:
        \begin{itemize}
            \item Edit it directly (advanced)
            \item Copy/paste from examples
            \item Share your full project as a single JSON
        \end{itemize}

        \vspace{0.15cm}
        \textbf{The \texttt{libraries.txt} file:}

        \vspace{0.1cm}
        Lists Arduino libraries your sketch needs:
        \begin{itemize}
            \item One library per line
            \item Example: \texttt{DHT sensor library}
            \item Example: \texttt{PubSubClient}
            \item Wokwi fetches them automatically
        \end{itemize}
    \end{columns}
\end{frame}

\begin{frame}{Wokwi Documentation and Resources}
    \begin{columns}
        \column{0.5\textwidth}
        \textbf{Official docs (bookmark these):}
        \begin{itemize}
            \item Docs home: \url{https://docs.wokwi.com/}
            \item ESP32 guide: \url{https://docs.wokwi.com/guides/esp32}
            \item Component list: \url{https://docs.wokwi.com/parts/wokwi-dht22}
            \item WiFi simulation: \url{https://docs.wokwi.com/guides/esp32-wifi}
        \end{itemize}

        \vspace{0.2cm}
        \textbf{Example projects (learn by reading):}
        \begin{itemize}
            \item ESP32 + DHT22: \url{https://wokwi.com/projects/322410731508073042}
            \item ESP32 + MQTT: \url{https://wokwi.com/projects/322524997423727188}
        \end{itemize}

        \column{0.5\textwidth}
        \fcolorbox{aiborder}{aibg}{
        \begin{minipage}{0.9\textwidth}
            \vspace{0.1cm}
            \textbf{Using AI with Wokwi:}
            \small

            \vspace{0.1cm}
            You can ask ChatGPT / Copilot to generate Arduino code for Wokwi. But always:
            \begin{enumerate}
                \item Verify pin numbers match your wiring
                \item Check library names are correct
                \item Test edge cases (sensor failure, WiFi drop)
                \item Read the Wokwi docs for simulator-specific details (e.g., \texttt{Wokwi-GUEST} WiFi SSID)
            \end{enumerate}
            \vspace{0.1cm}
        \end{minipage}
        }

        \vspace{0.15cm}
        \fcolorbox{warnborder}{warnbg}{
        \begin{minipage}{0.9\textwidth}
            \vspace{0.1cm}
            \small
            \textbf{Tip:} Open an example project, click ``Duplicate'', then modify it. Faster than starting from scratch.
            \vspace{0.1cm}
        \end{minipage}
        }
    \end{columns}
\end{frame}

%==============================================================================
\section{Why Simulate First?}
%==============================================================================

\begin{frame}[shrink=5]{Why Simulate Before Building?}
    \begin{columns}
        \column{0.55\textwidth}
        \textbf{Professional firmware workflow:}
        \begin{enumerate}
            \item Write code on simulator
            \item Test all logic paths
            \item Flash to real hardware
            \item Debug only hardware-specific issues
        \end{enumerate}

        \vspace{0.2cm}
        \textbf{What you gain today:}
        \begin{itemize}
            \item Working firmware \textbf{before} hardware arrives (Mar 29)
            \item Tested sensor reading logic
            \item Tested MQTT publish code
            \item Tested WiFi connection handling
        \end{itemize}

        \column{0.45\textwidth}
        \fcolorbox{aiborder}{aibg}{
        \begin{minipage}{0.9\textwidth}
            \vspace{0.1cm}
            \textbf{AI era context:}

            \vspace{0.1cm}
            \small
            AI can generate ESP32 code. But:
            \begin{itemize}
                \item Does it handle WiFi reconnect?
                \item Does it manage sensor errors?
                \item Does it respect power budget?
                \item Does it match YOUR BOM?
            \end{itemize}

            \vspace{0.1cm}
            Today: write, test, and \textbf{understand} every line.
            \vspace{0.1cm}
        \end{minipage}
        }

        \vspace{0.2cm}
        \begin{alertblock}{Goal}
            \small
            Leave today with firmware you can flash on Mar 29.
        \end{alertblock}
    \end{columns}
\end{frame}

%==============================================================================
\section{Wokwi Setup}
%==============================================================================

\begin{frame}{Step 0: Open Wokwi (2 min)}
    \begin{enumerate}
        \item Open \url{https://wokwi.com/} in your browser
        \item Click \textbf{``Start from Scratch''} $\rightarrow$ \textbf{ESP32}
        \item You should see:
        \begin{itemize}
            \item \textbf{Left panel:} Code editor (Arduino C++)
            \item \textbf{Right panel:} Virtual ESP32 board
            \item \textbf{Green play button:} Run simulation
        \end{itemize}
        \item Click the \textbf{``+''} button on the board to add components
    \end{enumerate}

    \vspace{0.3cm}
    \fcolorbox{stepborder}{stepbg}{
    \begin{minipage}{0.95\textwidth}
        \small
        Works immediately, on any browser (Chrome recommended).\\
        Create a free account if you want to save projects.
    \end{minipage}
    }
\end{frame}

%==============================================================================
\section{Exercise 1: Sensor Reading}
%==============================================================================

\begin{frame}[fragile]{Exercise 1: Read a Virtual DHT22 (15 min)}
    \textbf{Add a DHT22 sensor to your ESP32 board:}
    \begin{enumerate}
        \item Click \textbf{``+''} on the board $\rightarrow$ search \textbf{``DHT22''}
        \item Wire: VCC $\rightarrow$ 3.3V, GND $\rightarrow$ GND, Data $\rightarrow$ GPIO 15
    \end{enumerate}

    \vspace{0.1cm}
    \begin{lstlisting}[language=C]
#include <DHT.h>
#define DHTPIN 15
#define DHTTYPE DHT22
DHT dht(DHTPIN, DHTTYPE);

void setup() {
  Serial.begin(115200);
  dht.begin();
  Serial.println("DHT22 initialized");
}

void loop() {
  float temp = dht.readTemperature();
  float hum = dht.readHumidity();
  if (isnan(temp) || isnan(hum)) {
    Serial.println("ERROR: Failed to read DHT22!");
    return;  // Handle sensor failure gracefully
  }
  Serial.printf("Temp: %.1f C | Humidity: %.1f %%\n", temp, hum);
  delay(2000);  // DHT22 min interval = 2 seconds
}
    \end{lstlisting}
\end{frame}

\begin{frame}[shrink=5]{Exercise 1: What to Verify}
    \textbf{Run the simulation and check:}

    \vspace{0.2cm}
    \begin{columns}
        \column{0.5\textwidth}
        \begin{enumerate}
            \item Open Serial Monitor (bottom panel)
            \item You should see temperature and humidity readings
            \item \textbf{Click on the DHT22} in the simulator to change values
            \item Test edge cases:
            \begin{itemize}
                \item Set temp to $-40$\textdegree{}C (sensor minimum)
                \item Set temp to $+80$\textdegree{}C (sensor maximum)
                \item Disconnect the data wire
            \end{itemize}
        \end{enumerate}

        \column{0.5\textwidth}
        \fcolorbox{warnborder}{warnbg}{
        \begin{minipage}{0.9\textwidth}
            \vspace{0.1cm}
            \textbf{Engineering questions:}
            \small
            \begin{itemize}
                \item What happens when the sensor fails?
                \item Does your code crash or handle it?
                \item What's the minimum read interval?
                \item How does this match the DHT22 datasheet?
            \end{itemize}
            \vspace{0.1cm}
        \end{minipage}
        }

        \vspace{0.15cm}
        \begin{block}{Connect to your BOM}
            \small
            If your project uses DHT22: this is your firmware.\\
            If not: swap for your sensor (search Wokwi library).
        \end{block}
    \end{columns}
\end{frame}

%==============================================================================
\section{Exercise 2: WiFi + MQTT}
%==============================================================================

\begin{frame}[fragile]{Exercise 2: WiFi + MQTT Publish (25 min)}
    \textbf{Connect your ESP32 to WiFi and publish sensor data via MQTT:}

    \vspace{0.1cm}
    \begin{lstlisting}[language=C]
#include <WiFi.h>
#include <PubSubClient.h>
#include <DHT.h>

#define DHTPIN 15
#define DHTTYPE DHT22
#define WIFI_SSID "Wokwi-GUEST"  // Wokwi virtual WiFi
#define WIFI_PASS ""
#define MQTT_BROKER "broker.hivemq.com"  // Free public broker
#define MQTT_PORT 1883

WiFiClient espClient;
PubSubClient mqtt(espClient);
DHT dht(DHTPIN, DHTTYPE);

void connectWiFi() {
  Serial.print("Connecting to WiFi...");
  WiFi.begin(WIFI_SSID, WIFI_PASS);
  while (WiFi.status() != WL_CONNECTED) {
    delay(500); Serial.print(".");
  }
  Serial.println(" Connected! IP: " + WiFi.localIP().toString());
}
    \end{lstlisting}
\end{frame}

\begin{frame}[fragile]{Exercise 2: MQTT Publish (continued)}
    \begin{lstlisting}[language=C]
void connectMQTT() {
  while (!mqtt.connected()) {
    Serial.print("Connecting to MQTT...");
    String clientId = "esp32-" + String(random(0xffff), HEX);
    if (mqtt.connect(clientId.c_str())) {
      Serial.println(" Connected!");
    } else {
      Serial.printf(" Failed (rc=%d). Retrying in 5s...\n",
                     mqtt.state());
      delay(5000);
    }
  }
}

void setup() {
  Serial.begin(115200);
  dht.begin();
  connectWiFi();
  mqtt.setServer(MQTT_BROKER, MQTT_PORT);
}

void loop() {
  if (!mqtt.connected()) connectMQTT();
  mqtt.loop();
  float t = dht.readTemperature();
  float h = dht.readHumidity();
  if (!isnan(t) && !isnan(h)) {
    String payload = "{\"temp\":" + String(t,1) +
                     ",\"hum\":" + String(h,1) + "}";
    mqtt.publish("eece5155/YOUR_NAME/sensors", payload.c_str());
    Serial.println("Published: " + payload);
  }
  delay(10000);  // Publish every 10 seconds
}
    \end{lstlisting}
\end{frame}

\begin{frame}[shrink=5]{Exercise 2: What to Verify}
    \begin{columns}
        \column{0.55\textwidth}
        \textbf{Test your MQTT connection:}
        \begin{enumerate}
            \item Run the simulation in Wokwi
            \item Open \url{https://www.hivemq.com/demos/websocket-client/}
            \item Subscribe to: \texttt{eece5155/YOUR\_NAME/sensors}
            \item You should see JSON payloads arriving
            \item Test resilience:
            \begin{itemize}
                \item What if WiFi drops? (Wokwi simulates this)
                \item What if broker is down?
                \item What if sensor returns NaN?
            \end{itemize}
        \end{enumerate}

        \vspace{0.1cm}
        \textbf{Customize for your project:}
        \begin{itemize}
            \item Change topic to match your MVP design
            \item Add your specific sensor type
            \item Adjust publish interval to your PES
        \end{itemize}

        \column{0.45\textwidth}
        \fcolorbox{aiborder}{aibg}{
        \begin{minipage}{0.9\textwidth}
            \vspace{0.1cm}
            \small
            \textbf{AI check:}

            \vspace{0.1cm}
            AI-generated MQTT code often:
            \begin{itemize}
                \item Forgets reconnection logic
                \item Hardcodes credentials
                \item Ignores sensor error handling
                \item Uses blocking delays
            \end{itemize}

            \vspace{0.1cm}
            Your code handles all of these.
            \vspace{0.1cm}
        \end{minipage}
        }

        \vspace{0.1cm}
        \begin{alertblock}{Deliverable}
            \small
            Save your Wokwi project link.\\
            This firmware goes to real ESP32 on Mar 29.
        \end{alertblock}
    \end{columns}
\end{frame}

%==============================================================================
\section{Exercise 3: JSON + Multiple Sensors}
%==============================================================================

\begin{frame}[fragile,shrink=5]{Exercise 3: Structured JSON + Multiple Sensors (20 min)}
    \textbf{Real deployments need structured, parseable data:}

    \vspace{0.1cm}
    \begin{lstlisting}[language=C]
// Add a photoresistor (LDR) to GPIO 34 (analog input)
#define LDR_PIN 34

void publishSensorData() {
  float temp = dht.readTemperature();
  float hum = dht.readHumidity();
  int light = analogRead(LDR_PIN);  // 0-4095 (12-bit ADC)
  float voltage = light * (3.3 / 4095.0);

  // Structured JSON with metadata
  String json = "{";
  json += "\"device_id\":\"esp32-node-01\",";
  json += "\"timestamp\":" + String(millis()) + ",";
  json += "\"sensors\":{";
  json += "\"temperature\":{\"value\":" + String(temp,1)
        + ",\"unit\":\"C\"},";
  json += "\"humidity\":{\"value\":" + String(hum,1)
        + ",\"unit\":\"%\"},";
  json += "\"light\":{\"raw\":" + String(light)
        + ",\"voltage\":" + String(voltage,2) + "}";
  json += "}}";

  mqtt.publish("eece5155/YOUR_NAME/data", json.c_str());
}
    \end{lstlisting}

    \vspace{0.1cm}
    \fcolorbox{stepborder}{stepbg}{
    \begin{minipage}{0.95\textwidth}
        \small
        \textbf{Why structured JSON?} Your Node-RED flows (Session 2) and InfluxDB (cloud hands-on) parse this format. Consistent schema = reliable pipeline.
    \end{minipage}
    }
\end{frame}

\begin{frame}[shrink=5]{Exercise 3: Adapt to YOUR Project}
    \begin{columns}
        \column{0.5\textwidth}
        \textbf{Wokwi has these virtual sensors:}

        \vspace{0.15cm}
        \footnotesize
        \begin{tabular}{ll}
            \toprule
            \textbf{Sensor} & \textbf{Use case} \\
            \midrule
            DHT22 & Temperature + humidity \\
            DS18B20 & Waterproof temperature \\
            HC-SR04 & Ultrasonic distance \\
            PIR & Motion detection \\
            Potentiometer & Analog input (any) \\
            LDR & Light level \\
            MPU6050 & Accelerometer + gyro \\
            Servo & Actuator output \\
            LCD 16x2 & Local display \\
            NeoPixel & LED status indicators \\
            \bottomrule
        \end{tabular}
        \normalsize

        \column{0.5\textwidth}
        \textbf{Map to your BOM:}

        \vspace{0.15cm}
        \begin{enumerate}
            \item Find the closest Wokwi sensor to what you ordered
            \item Write the read function
            \item Publish via MQTT with your topic schema
            \item Test edge cases (min/max/error)
        \end{enumerate}

        \vspace{0.15cm}
        \fcolorbox{warnborder}{warnbg}{
        \begin{minipage}{0.9\textwidth}
            \vspace{0.1cm}
            \small
            \textbf{Sensor missing from Wokwi?}

            \vspace{0.05cm}
            Use a potentiometer as analog proxy.\\
            Map 0--4095 to your sensor's range.\\
            Example: soil moisture 0--100\%.
            \vspace{0.1cm}
        \end{minipage}
        }
    \end{columns}
\end{frame}

%==============================================================================
\section{Wrap-Up}
%==============================================================================

\begin{frame}{Deliverables + Next Steps}
    \textbf{Before you leave today:}

    \vspace{0.2cm}
    \begin{enumerate}
        \item \textbf{Save your Wokwi project} (create account $\rightarrow$ save)
        \item \textbf{Post your Wokwi link} on Canvas discussion board
        \item \textbf{Verify:} Your MQTT messages appear on HiveMQ web client
    \end{enumerate}

    \vspace{0.3cm}
    \textbf{This week's sessions build on each other:}

    \vspace{0.15cm}
    \begin{tabular}{lll}
        \toprule
        \textbf{Session} & \textbf{What you build} & \textbf{Connects to} \\
        \midrule
        \textbf{Today (Mon 23)} & ESP32 firmware + MQTT publish & Device side \\
        \textbf{Wed 25} & Node-RED flows + dashboards & Automation side \\
        \textbf{Thu 26} & IoT dataset analysis & Data quality skills \\
        \midrule
        \textbf{Mar 29+} & Real hardware in lab & Everything connects \\
        \bottomrule
    \end{tabular}

    \vspace{0.3cm}
    \fcolorbox{stepborder}{stepbg}{
    \begin{minipage}{0.95\textwidth}
        \small
        \textbf{Portfolio note:} Include your Wokwi project link in your final report as evidence of ``firmware development and simulation testing before hardware deployment.''
    \end{minipage}
    }
\end{frame}

\end{document}
